\documentclass[12pt]{article}
\usepackage{amsmath}
\usepackage{amssymb}
\usepackage{physics}
\usepackage{hyperref}

\newcommand{\ds}{\displaystyle}

\title{The Relativistic Maxwell Equations}
\author{}
\date{}

\begin{document}
\maketitle


% ============================================================
\section{Introduction}
% ============================================================
In this segment, we derive the transformation of Maxwell's equations under a Lorentz boost in
an arbitrary direction.

Modern treatments almost invariably rely on the language of four-vectors and tensors, through
which the relativistic covariance of Maxwell's equations follows naturally. While this
formalism is undoubtedly the most elegant framework for relativistic electrodynamics, it often
conceals the multivariable calculus underlying the transformation of the familiar differential
operators.

Our approach follows a different path. Inspired by Einstein's original derivation, we work
entirely within the framework of classical multivariable calculus, extending the analysis from
boosts along the coordinate axes to a Lorentz boost in an arbitrary direction. Although this
approach is considerably longer than the tensor derivation, it exposes every intermediate step
and makes explicit how the gradient, divergence, curl, and time derivative transform between
inertial reference frames. Understanding these transformations in detail provides valuable
intuition for the geometric ideas that the tensor formalism later encapsulates.

We begin by introducing the multivariable calculus identities required throughout the chapter.
We then derive the transformation laws for the differential operators under a Lorentz
transformation. Finally, these results are combined to obtain the relativistic form of
Maxwell's equations in a charge-free vacuum.

\subsection{Definitions and Useful Formulas from Multivariate Calculus}

We will consider a stationary system of reference where the coordinates of an event are
represented by the four-vector $(x_1, x_2, x_3, t)$, and a moving frame where the same event
has coordinates $x'_1, x'_2, x'_3, t'$. The fourth coordinate in both four-vectors has units
of time, not space, as the mathematics of the Maxwell equations derivation keeps the notation considerably 
lighter throughout the derivation.

A vector field $\vb{A}$ is a three-dimensional valued function acting on a four-vector,
\[
  \begin{gathered}
    \vb{A}(x_1, x_2, x_3, t) \\
    \vb{A}: \mathbb{R}^4 \to \mathbb{R}^3
  \end{gathered}
\]

A few differential operators of $\vb{A}$ are extensively used in the derivation below and
deserve to be explicitly defined to minimize confusion when the notation will get a little
heavy, as it often does in the derivation of the relativistic Maxwell equations.

\begin{itemize}
  \item The Jacobian on the spatial dimensions and time:
  \[
    \nabla \vb{A}_{x,t} =
    \begin{pmatrix}
      \pdv{A_1}{x_1} & \pdv{A_1}{x_2} & \pdv{A_1}{x_3} & \pdv{A_1}{t} \\[6pt]
      \pdv{A_2}{x_1} & \pdv{A_2}{x_2} & \pdv{A_2}{x_3} & \pdv{A_2}{t} \\[6pt]
      \pdv{A_3}{x_1} & \pdv{A_3}{x_2} & \pdv{A_3}{x_3} & \pdv{A_3}{t}
    \end{pmatrix}
  \]

  \item The Jacobian on the spatial dimensions:
  \[
    \nabla \vb{A}_{x} =
    \begin{pmatrix}
      \pdv{A_1}{x_1} & \pdv{A_1}{x_2} & \pdv{A_1}{x_3} \\[6pt]
      \pdv{A_2}{x_1} & \pdv{A_2}{x_2} & \pdv{A_2}{x_3} \\[6pt]
      \pdv{A_3}{x_1} & \pdv{A_3}{x_2} & \pdv{A_3}{x_3}
    \end{pmatrix}
  \]

\item The product of the spatial Jacobian with a vector $\vb{v}$, typically the
      relative velocity between two inertial frames:
\[
  \nabla_x\vb{A}\,\vb{v}
  =
  \begin{pmatrix}
    \pdv{A_1}{x_1}v_1+\pdv{A_1}{x_2}v_2+\pdv{A_1}{x_3}v_3 \\[6pt]
    \pdv{A_2}{x_1}v_1+\pdv{A_2}{x_2}v_2+\pdv{A_2}{x_3}v_3 \\[6pt]
    \pdv{A_3}{x_1}v_1+\pdv{A_3}{x_2}v_2+\pdv{A_3}{x_3}v_3
  \end{pmatrix}.
\]

  \item The divergence of a vector field $\vb{A}$:
  \[
    \div \vb{A} = \pdv{A_1}{x_1} + \pdv{A_2}{x_2} + \pdv{A_3}{x_3}
  \]

  \item And finally the curl of a vector field $\vb{A}$:
  \[
    \curl \vb{A} =
      \left(\pdv{A_3}{x_2} - \pdv{A_2}{x_3}\right)\vb{e}_1
      - \left(\pdv{A_3}{x_1} - \pdv{A_1}{x_3}\right)\vb{e}_2
      + \left(\pdv{A_2}{x_1} - \pdv{A_1}{x_2}\right)\vb{e}_3
  \]
  where $\vb{e}_1$, $\vb{e}_2$ and $\vb{e}_3$ are unit vectors along the three spatial axes.
\end{itemize}

A few identities from calculus are key to derive the relativistic Maxwell equations, and they
are listed below:

\begin{itemize}
  \item The divergence and the curl of the cross product of a constant vector $\vb{v}$ and a
        field $\vb{A}$:
  \begin{equation}
    \begin{aligned}
      \curl(\vb{v} \cross \vb{A}) &= -\nabla_x \vb{A} \cdot \vb{v} + (\div \vb{A})\,\vb{v} \\
      \div(\vb{v} \cross \vb{A})  &= -\ip{\vb{v}}{\curl \vb{A}}
    \end{aligned}
    \label{eq:divergence_curl_cross_product}
  \end{equation}

  \item The curl of the triple product of a constant vector $\vb{v}$ with a field $\vb{A}$,
        which has two different expressions, the latter derived with the use of the equations
        above:
  \begin{equation}
    \begin{aligned}
      \curl\bigl(\vb{v} \cross (\vb{v} \cross \vb{A})\bigr)
        &= \curl\bigl(\ip{\vb{v}}{\vb{A}}\vb{v} - v^2\vb{A}\bigr) \\[4pt]
      \curl\bigl(\vb{v} \cross (\vb{v} \cross \vb{A})\bigr)
        &= -\nabla_x(\vb{v} \cross \vb{A}) \cdot \vb{v}
           + \div(\vb{v} \cross \vb{A})\,\vb{v} \\
        &= -\vb{v} \cross (\nabla_x \vb{A} \cdot \vb{v})
           - \ip{\vb{v}}{\curl \vb{A}}\,\vb{v}
    \end{aligned}
    \label{eq:curl_triple_product}
  \end{equation}
\end{itemize}

% ============================================================
\section{Transformation of Differential Operators}
% ============================================================

Before deriving the relativistic Maxwell equations, we first determine how the differential operators transform under a 
Lorentz transformation. We remind the reader of the definition,  

\[
\gamma=\frac{1}{\sqrt{1-\frac{v^2}{c^2}}}.
\]

We consider a field $\vb{A}$ in $F$,

\[
  \begin{gathered}
    \vb{A} = A(x_1, x_2, x_3, t) \\
    \vb{A}: \mathbb{R}^4 \to \mathbb{R}^3
  \end{gathered}
\]

which we compose with the inverse of the Lorentz transformation from $F'$ to $F$ (see
\textit{The Transformation of Time and Space}) obtaining a field $\vb{A}'$ defined on the
space $F'$,

\begin{align}
  \vb{A}' &= \vb{A} \circ \boldsymbol{\Lambda}(-\vb{v}) \\
  \vb{A}' &= \vb{A}'(\vb{x}', t') = \vb{A}\bigl(\boldsymbol{\Lambda}(-\vb{v})(\vb{x}', t')\bigr) \nonumber \\
  \vb{A}' &: \mathbb{R}^4 \to \mathbb{R}^3 \nonumber
\end{align}

Since the Maxwell equations are in $F$ coordinates, we need to understand how the partial
derivatives of the field $\vb{A}'$ transform under the Lorentz transformation, which we are
able to compute from the Jacobian chain rule,

\[
  \nabla_{x',t'}\vb{A}' = \nabla_{x,t}\vb{A} \cdot \nabla_{x',t'}\bigl(\boldsymbol{\Lambda}(-\vb{v})(\vb{x}',t')\bigr)
\]

where $\nabla_{x',t'}\vb{A}'$ and $\nabla_{x,t}\vb{A}$ are the Jacobians of $\vb{A}'$ and
$\vb{A}$ respectively, computed in their own coordinate systems.

The Lorentz transformation is linear, hence the Jacobian is $\Lambda$, and the chain rule
becomes, in coordinate form,

\[
  \begin{pmatrix}
    \ds\pdv{A'_1}{x'_1} & \ds\pdv{A'_1}{x'_2} & \ds\pdv{A'_1}{x'_3} & \ds\pdv{A'_1}{t'} \\[8pt]
    \ds\pdv{A'_2}{x'_1} & \ds\pdv{A'_2}{x'_2} & \ds\pdv{A'_2}{x'_3} & \ds\pdv{A'_2}{t'} \\[8pt]
    \ds\pdv{A'_3}{x'_1} & \ds\pdv{A'_3}{x'_2} & \ds\pdv{A'_3}{x'_3} & \ds\pdv{A'_3}{t'}
  \end{pmatrix}
  =
  \begin{pmatrix}
    \ds\pdv{A_1}{x_1} & \ds\pdv{A_1}{x_2} & \ds\pdv{A_1}{x_3} & \ds\pdv{A_1}{t} \\[8pt]
    \ds\pdv{A_2}{x_1} & \ds\pdv{A_2}{x_2} & \ds\pdv{A_2}{x_3} & \ds\pdv{A_2}{t} \\[8pt]
    \ds\pdv{A_3}{x_1} & \ds\pdv{A_3}{x_2} & \ds\pdv{A_3}{x_3} & \ds\pdv{A_3}{t}
  \end{pmatrix}
  \cdot\, \Lambda(-\vb{v})
\]

To get to the partial derivatives of $\vb{A}$ we need to invert the equation above,
\[
  \begin{pmatrix}
    \ds\pdv{A'_1}{x'_1} & \ds\pdv{A'_1}{x'_2} & \ds\pdv{A'_1}{x'_3} & \ds\pdv{A'_1}{t'} \\[8pt]
    \ds\pdv{A'_2}{x'_1} & \ds\pdv{A'_2}{x'_2} & \ds\pdv{A'_2}{x'_3} & \ds\pdv{A'_2}{t'} \\[8pt]
    \ds\pdv{A'_3}{x'_1} & \ds\pdv{A'_3}{x'_2} & \ds\pdv{A'_3}{x'_3} & \ds\pdv{A'_3}{t'}
  \end{pmatrix}
  \cdot\,\Lambda(\vb{v})
  =
  \begin{pmatrix}
    \ds\pdv{A_1}{x_1} & \ds\pdv{A_1}{x_2} & \ds\pdv{A_1}{x_3} & \ds\pdv{A_1}{t} \\[8pt]
    \ds\pdv{A_2}{x_1} & \ds\pdv{A_2}{x_2} & \ds\pdv{A_2}{x_3} & \ds\pdv{A_2}{t} \\[8pt]
    \ds\pdv{A_3}{x_1} & \ds\pdv{A_3}{x_2} & \ds\pdv{A_3}{x_3} & \ds\pdv{A_3}{t}
  \end{pmatrix}
\]

The bra-ket formalism helps us here once again to compute the product of the Jacobian
matrices. Any row $j$ of $\nabla_{x',t'}\vb{A}' \cdot \Lambda(\vb{v})$ is the vector
\[
  \begin{aligned}
    &\bra{\ds\pdv{A'_j}{x'_1},\, \ds\pdv{A'_j}{x'_2},\, \ds\pdv{A'_j}{x'_3},\, \ds\pdv{A'_j}{t'}}
    \left[
      \frac{\gamma}{v^2}\ket{\vb{v},0}\bra{\vb{v},0}
      - \gamma\ket{\vb{v},0}\bra{\vb{0},1} \right.\\
    &\left.
      + \left(\vb{I} - \frac{1}{v^2}\ket{\vb{v},0}\bra{\vb{v},0}\right)
      - \frac{\gamma}{c^2}\ket{\vb{0},1}\bra{\vb{v},0}
      + \gamma\ket{\vb{0},1}\bra{\vb{0},1}
    \right]
  \end{aligned}
\]
where $\vb{I}$ is the identity operator, and reads,
\[
  \begin{aligned}
    &\frac{\gamma}{v^2}\ip{\nabla_{x'}\vb{A}'_j}{\vb{v}}\bra{\vb{v},0}
    - \gamma\ip{\nabla_{x'}\vb{A}'_j}{\vb{v}}\bra{\vb{0},1}
    + \bra{\nabla_{x'}\vb{A}'_j,\,0} \\
    &- \frac{1}{v^2}\ip{\nabla_{x'}\vb{A}'_j}{\vb{v}}\bra{\vb{v},0}
    - \frac{\gamma}{c^2}\ds\pdv{A'_j}{t'}\bra{\vb{v},0}
    + \gamma\ds\pdv{A'_j}{t'}\bra{\vb{0},1}
  \end{aligned}
\]

The spatial gradient of any component of the field $\vb{A}$ computed at
$(\vb{x},t) = \Lambda(-\vb{v})(\vb{x}',t')$ is,
\begin{equation}
  \nabla_x \vb{A}_j = \nabla_{x'}\vb{A}'_j
    + \frac{\gamma-1}{v^2}\ip{\nabla_{x'}\vb{A}'_j}{\vb{v}}\bra{\vb{v}}
    - \frac{\gamma}{c^2}\ds\pdv{A'_j}{t'}\bra{\vb{v}}
  \label{eq:gradient_in_F}
\end{equation}
and the time derivative is
\begin{equation}
  \pdv{A_j}{t} = -\gamma\ip{\nabla_{x'}\vb{A}'_j}{\vb{v}} + \gamma\pdv{A'_j}{t'}
  \label{eq:time_derivative_in_F}
\end{equation}
where,

\begin{align}
  \nabla_{x'}\vb{A}'_j &= \bra{\ds\pdv{A'_j}{x'_1},\, \ds\pdv{A'_j}{x'_2},\, \ds\pdv{A'_j}{x'_3}} \\
  \nabla_x \vb{A}_j    &= \bra{\ds\pdv{A_j}{x_1},\, \ds\pdv{A_j}{x_2},\, \ds\pdv{A_j}{x_3}}
\end{align}

Equations \eqref{eq:gradient_in_F} and \eqref{eq:time_derivative_in_F} provide the transformation law for the first-order 
partial derivatives of any vector field. 
We now use these expressions to derive the transformation laws for the differential operators appearing in Maxwell's equations.
\subsection{Divergence in $F$ and $F'$}

We now use \eqref{eq:gradient_in_F} to compute the divergence of the field $\vb{A}$ in $F$,
\begin{align}
  \pdv{A_1}{x_1} &= \pdv{A'_1}{x'_1}
    + \frac{\gamma-1}{v^2}\ip{\nabla_{x'}\vb{A}'_1}{\vb{v}}v_1
    - \frac{\gamma}{c^2}\pdv{A'_1}{t'}v_1 \\
  \pdv{A_2}{x_2} &= \pdv{A'_2}{x'_2}
    + \frac{\gamma-1}{v^2}\ip{\nabla_{x'}\vb{A}'_2}{\vb{v}}v_2
    - \frac{\gamma}{c^2}\pdv{A'_2}{t'}v_2 \\
  \pdv{A_3}{x_3} &= \pdv{A'_3}{x'_3}
    + \frac{\gamma-1}{v^2}\ip{\nabla_{x'}\vb{A}'_3}{\vb{v}}v_3
    - \frac{\gamma}{c^2}\pdv{A'_3}{t'}v_3
\end{align}

To derive an expression for the field divergence, we expand the scalar products in the
equation above and group the terms on the partial derivative along each coordinate. We will
work out the details for the first equation, as the others are identically derived by just
changing the indexes.
\[
  \begin{aligned}
    \pdv{A_1}{x_1} &= \pdv{A'_1}{x'_1}
      + \frac{\gamma-1}{v^2}\ip{\nabla_{x'}\vb{A}'_1}{\vb{v}}v_1
      - \frac{\gamma}{c^2}\pdv{A'_1}{t'}v_1 \\
    &= \pdv{A'_1}{x'_1}
      + \frac{\gamma-1}{v^2}\!\left(
          \pdv{A'_1}{x'_1}v_1 + \pdv{A'_1}{x'_2}v_2 + \pdv{A'_1}{x'_3}v_3
        \right)v_1
      - \frac{\gamma}{c^2}\pdv{A'_1}{t'}v_1 \\
    &= \left(1 + \frac{\gamma-1}{v^2}v_1^2\right)\pdv{A'_1}{x'_1}
      + \frac{\gamma-1}{v^2}\!\left(
          \pdv{(A'_1 v_1)}{x'_2}v_2 + \pdv{(A'_1 v_1)}{x'_3}v_3
        \right)
      - \frac{\gamma}{c^2}\pdv{(A'_1 v_1)}{t'}
  \end{aligned}
\]

The divergence of the field in $F$ is,
\[
  \begin{aligned}
    &\pdv{A_1}{x_1} + \pdv{A_2}{x_2} + \pdv{A_3}{x_3} \\
    &\quad= \left(1 + \frac{\gamma-1}{v^2}v_1^2\right)\pdv{A'_1}{x'_1}
      + \frac{\gamma-1}{v^2}\!\left(\pdv{(A'_1 v_1)}{x'_2}v_2 + \pdv{(A'_1 v_1)}{x'_3}v_3\right)
      - \frac{\gamma}{c^2}\pdv{(A'_1 v_1)}{t'} \\
    &\quad+ \left(1 + \frac{\gamma-1}{v^2}v_2^2\right)\pdv{A'_2}{x'_2}
      + \frac{\gamma-1}{v^2}\!\left(\pdv{(A'_2 v_2)}{x'_1}v_1 + \pdv{(A'_2 v_2)}{x'_3}v_3\right)
      - \frac{\gamma}{c^2}\pdv{(A'_2 v_2)}{t'} \\
    &\quad+ \left(1 + \frac{\gamma-1}{v^2}v_3^2\right)\pdv{A'_3}{x'_3}
      + \frac{\gamma-1}{v^2}\!\left(\pdv{(A'_3 v_3)}{x'_1}v_1 + \pdv{(A'_3 v_3)}{x'_2}v_2\right)
      - \frac{\gamma}{c^2}\pdv{(A'_3 v_3)}{t'}
  \end{aligned}
\]

By grouping the terms by the components of the field $\vb{A}$, we obtain the divergence of
the field $\vb{A}$ in $F$,
\[
  \begin{aligned}
    &\pdv{A_1}{x_1} + \pdv{A_2}{x_2} + \pdv{A_3}{x_3} \\
    &= \pdv{A'_1}{x'_1} + \pdv{A'_2}{x'_2} + \pdv{A'_3}{x'_3} \\
    &\quad + \frac{\gamma-1}{v^2}\!\left(\pdv{A'_1}{x'_1}v_1 + \pdv{A'_1}{x'_2}v_2 + \pdv{A'_1}{x'_3}v_3\right)v_1 \\
    &\quad + \frac{\gamma-1}{v^2}\!\left(\pdv{A'_2}{x'_1}v_1 + \pdv{A'_2}{x'_2}v_2 + \pdv{A'_2}{x'_3}v_3\right)v_2 \\
    &\quad + \frac{\gamma-1}{v^2}\!\left(\pdv{A'_3}{x'_1}v_1 + \pdv{A'_3}{x'_2}v_2 + \pdv{A'_3}{x'_3}v_3\right)v_3 \\
    &\quad - \frac{\gamma}{c^2}\pdv{(A'_1 v_1 + A'_2 v_2 + A'_3 v_3)}{t'}
  \end{aligned}
\]

We get the compact equation for the transformation of the divergence of $\vb{A}$ to $\vb{A}'$,
\[
  \div \vb{A}
  =
  \div \vb{A}'
  + \frac{\gamma-1}{v^2}
    \ip{\nabla_{x'}\vb{A}'\,\vb v}{\vb v}
  - \frac{\gamma}{c^2}
    \pdv{\ip{\vb{A}'}{\vb{v}}}{t'}
\]

If we consider a divergence-free field $\vb{A}$ in $F$, the equation above gives an important
condition on the divergence of $\vb{A}'$,

\begin{equation}
  \div \vb{A}'
  =
  -\frac{\gamma-1}{v^2}
    \ip{\nabla_{x'}\vb{A}'\,\vb v}{\vb v}
  + \frac{\gamma}{c^2}
    \pdv{\ip{\vb{A}'}{\vb{v}}}{t'}
  \label{eq:divergence_in_Fprime}
\end{equation}

\subsection{Curl in $F$ and $F'$}

The other operator needed for the Maxwell equations is the curl,
\[
  \begin{aligned}
    (\curl \vb{A})_1 &= \pdv{A_3}{x_2} - \pdv{A_2}{x_3} \\
      &= \pdv{A'_3}{x'_2}
        + \frac{\gamma-1}{v^2}\ip{\nabla_{x'}\vb{A}'_3}{\vb{v}}v_2
        - \frac{\gamma}{c^2}\pdv{A'_3}{t'}v_2 \\
      &\quad - \pdv{A'_2}{x'_3}
        - \frac{\gamma-1}{v^2}\ip{\nabla_{x'}\vb{A}'_2}{\vb{v}}v_3
        + \frac{\gamma}{c^2}\pdv{A'_2}{t'}v_3 \\[6pt]
   (\curl \vb{A})_2 &= -\pdv{A_3}{x_1} + \pdv{A_1}{x_3} \\
  &= -\pdv{A'_3}{x'_1}
    - \frac{\gamma-1}{v^2}
      \ip{\nabla_{x'}\vb{A}'_3}{\vb{v}}v_1
    + \frac{\gamma}{c^2}\pdv{A'_3}{t'}v_1 \\
  &\quad + \pdv{A'_1}{x'_3}
    + \frac{\gamma-1}{v^2}
      \ip{\nabla_{x'}\vb{A}'_1}{\vb{v}}v_3
    - \frac{\gamma}{c^2}\pdv{A'_1}{t'}v_3 \\[6pt]
    (\curl \vb{A})_3 &= \pdv{A_2}{x_1} - \pdv{A_1}{x_2} \\
      &= \pdv{A'_2}{x'_1}
        + \frac{\gamma-1}{v^2}\ip{\nabla_{x'}\vb{A}'_2}{\vb{v}}v_1
        - \frac{\gamma}{c^2}\pdv{A'_2}{t'}v_1 \\
      &\quad - \pdv{A'_1}{x'_2}
        - \frac{\gamma-1}{v^2}\ip{\nabla_{x'}\vb{A}'_1}{\vb{v}}v_2
        + \frac{\gamma}{c^2}\pdv{A'_1}{t'}v_2
  \end{aligned}
\]

which is in vector form easily recognisable as,

\begin{equation}
  \curl \vb{A}
  =
  \curl \vb{A}'
  + \frac{\gamma-1}{v^2}\,
    \vb v \cross
    \left(
      \nabla_{x'}\vb A'\,\vb v
    \right)
  - \frac{\gamma}{c^2}
    \pdv{\vb v\cross\vb A'}{t'}
  \label{eq:curl_in_F}
\end{equation}

% ============================================================
\section{The Relativistic Maxwell Equations}
% ============================================================

We consider an electric field $\vb{E}$ and a magnetic field $\vb{B}$ in a region of the
system of reference $F$ free of charges and currents. The Maxwell equations in this case are

\begin{equation}
    \begin{aligned}
      \pdv{\vb{B}}{t} &= -\curl \vb{E} \\
      \pdv{\vb{E}}{t} &= c^2 \curl \vb{B} \\
      \div \vb{B}     &= 0 \\
      \div \vb{E}     &= 0
    \end{aligned}   
  \label{eq:maxwell_in_F}
\end{equation}

The fields $\vb{B}'$ and $\vb{E}'$, obtained by applying the composition of the Lorentz
transformation are,

\begin{align}
  \vb{B}' &= \vb{B}'(\vb{x}',t') = \vb{B}\bigl(\Lambda(-\vb{v})(\vb{x}',t')\bigr) \\
  \vb{E}' &= \vb{E}'(\vb{x}',t') = \vb{E}\bigl(\Lambda(-\vb{v})(\vb{x}',t')\bigr)
\end{align}

and fulfill the equations obtained by applying directly the Lorentz transformation to
equations \eqref{eq:maxwell_in_F} (see section \textbf{Transformation of Differential
Operators}),

\begin{equation}
  \begin{aligned}
    -\gamma\nabla_{x'}\vb{B}'\,\vb{v}
    + \gamma\pdv{\vb{B}'}{t'}
    &=
    -\curl \vb{E}'
    - \frac{\gamma-1}{v^2}\,
      \vb{v}\cross\left(\nabla_{x'}\vb{E}'\,\vb{v}\right)
    + \frac{\gamma}{c^2}
      \pdv{\vb{v}\cross\vb{E}'}{t'},
    \\
    -\frac{\gamma}{c^2}\nabla_{x'}\vb{E}'\,\vb{v}
    + \frac{\gamma}{c^2}\pdv{\vb{E}'}{t'}
    &=
    \curl \vb{B}'
    + \frac{\gamma-1}{v^2}\,
      \vb{v}\cross\left(\nabla_{x'}\vb{B}'\,\vb{v}\right)
    - \frac{\gamma}{c^2}
      \pdv{\vb{v}\cross\vb{B}'}{t'}.
  \end{aligned}
  \label{eq:maxwell_transformed}
\end{equation}

Since the equations at hand do not conform to the standard form of Maxwell's equations, the
fields denoted by $\vb{B}'$ and $\vb{E}'$ do not represent the transformed electric and
magnetic fields we seek. Substantial effort is required to manipulate these equations into the
correct form that yields the desired fields.

We outline below the steps required to derive the relativistic Maxwell equations in a vacuum,
starting with the equations above and manipulating them to obtain the correct form.

 
\subsection{Step 1: Divergence-Free Fields}

Taking the scalar product of equations \eqref{eq:maxwell_transformed} with the velocity
$\vb{v}$ gives

\begin{align}
  -\gamma\ip{\nabla_{x'}\vb{B}'\,\vb{v}}{\vb{v}}
    + \gamma\pdv{\ip{\vb{B}'}{\vb{v}}}{t'}
    &= -\ip{\curl \vb{E}'}{\vb{v}},
    \\
  -\frac{\gamma}{c^2}\ip{\nabla_{x'}\vb{E}'\,\vb{v}}{\vb{v}}
    + \frac{\gamma}{c^2}\pdv{\ip{\vb{E}'}{\vb{v}}}{t'}
    &= \ip{\curl \vb{B}'}{\vb{v}}.
\end{align}

Together with the divergence-free condition \eqref{eq:divergence_in_Fprime}, these equations
give useful expressions for the divergences of $\vb{B}'$ and $\vb{E}'$:

\begin{align}
  \div \vb{B}'
    &=
    \frac{\gamma-1}{\gamma v^2}
    \left(
      \pdv{\ip{\vb{B}'}{\vb{v}}}{t'}
      -
      \ip{\curl \vb{E}'}{\vb{v}}
    \right),
    \\
  \div \vb{E}'
    &=
    \frac{\gamma-1}{\gamma v^2}
    \left(
      \pdv{\ip{\vb{E}'}{\vb{v}}}{t'}
      +
      c^2\ip{\curl \vb{B}'}{\vb{v}}
    \right).
\end{align}

The terms $\nabla_{x'}\vb{B}'\,\vb{v}$ and
$\nabla_{x'}\vb{E}'\,\vb{v}$ on the left-hand side of
\eqref{eq:maxwell_transformed} can be substituted using
\eqref{eq:divergence_curl_cross_product}. Using also the expressions for the divergences of
$\vb{B}'$ and $\vb{E}'$ just obtained, we find
\begin{equation}
  \begin{aligned}
    \pdv{}{t'}\!\left(
        \gamma\,\vb{B}'
        -
        \frac{\gamma-1}{v^2}\ip{\vb{B}'}{\vb{v}}\vb{v}
        -
        \frac{\gamma}{c^2}\vb{v}\cross\vb{E}'
      \right)
    &=
      -\gamma\curl(\vb{v}\cross\vb{B}')
      -
      \curl\vb{E}'
      \\
    &\quad
      -
      \frac{\gamma-1}{v^2}
      \ip{\curl\vb{E}'}{\vb{v}}\vb{v}
      -
      \frac{\gamma-1}{v^2}\,
      \vb{v}\cross
      \left(
        \nabla_{x'}\vb{E}'\,\vb{v}
      \right),
    \\[6pt]
    \pdv{}{t'}\!\left(
        \gamma\,\vb{E}'
        -
        \frac{\gamma-1}{v^2}\ip{\vb{E}'}{\vb{v}}\vb{v}
        +
        \gamma\,\vb{v}\cross\vb{B}'
      \right)
    &=
      -\gamma\curl(\vb{v}\cross\vb{E}')
      +
      \curl(c^2\vb{B}')
      \\
    &\quad
      +
      \frac{\gamma-1}{v^2}
      \ip{\curl(c^2\vb{B}')}{\vb{v}}\vb{v}
      +
      \frac{\gamma-1}{v^2}\,
      \vb{v}\cross
      \left(
        \nabla_{x'}(c^2\vb{B}')\,\vb{v}
      \right).
  \end{aligned}
  \label{eq:divergence_free_fields}
\end{equation}

We are nearing the solution. It is now necessary to evaluate the curl of the triple
cross product.
 
\subsection{Step 2: The Curl of the Triple Product}

Equations \eqref{eq:divergence_free_fields} bring us closer to the solution. On the
left-hand side, we have the partial derivatives with respect to $t'$ of two fields. On the
right-hand side, several terms are already expressed as curls, but the following terms are
not:
\[
  -\frac{\gamma-1}{v^2}\ip{\curl \vb{E}'}{\vb{v}}\vb{v}
  -
  \frac{\gamma-1}{v^2}\,
  \vb{v}\cross
  \left(
    \nabla_{x'}\vb{E}'\,\vb{v}
  \right).
\]

The key to the solution is the curl of the triple cross product, given by
\eqref{eq:curl_triple_product}:

\[
  \begin{aligned}
    \curl\bigl(\vb{v}\cross(\vb{v}\cross\vb{E}')\bigr)
    &=
    \curl\bigl(
      \ip{\vb{v}}{\vb{E}'}\vb{v}
      -
      v^2\vb{E}'
    \bigr)
    \\
    &=
    -\vb{v}\cross
    \left(
      \nabla_{x'}\vb{E}'\,\vb{v}
    \right)
    -
    \ip{\vb{v}}{\curl\vb{E}'}\vb{v}.
  \end{aligned}
\]

Substitution into \eqref{eq:divergence_free_fields} finally yields equations having the
same differential-operator form as \eqref{eq:maxwell_in_F}, but now expressed in the
coordinates of $F'$:
\begin{equation}
  \begin{aligned}
    \pdv{}{t'}\!\left(
        \gamma\,\vb{B}'
        -
        (\gamma-1)
        \ip{\vb{B}'}{\frac{\vb{v}}{v}}
        \frac{\vb{v}}{v}
        -
        \frac{\gamma}{c^2}\vb{v}\cross\vb{E}'
      \right)
    &=
    -\curl\!\left(
        \gamma\,\vb{E}'
        -
        (\gamma-1)
        \ip{\vb{E}'}{\tfrac{\vb{v}}{v}}
        \tfrac{\vb{v}}{v}
        +
        \gamma\,\vb{v}\cross\vb{B}'
      \right),
    \\[6pt]
    \pdv{}{t'}\!\left(
        \gamma\,\vb{E}'
        -
        (\gamma-1)
        \ip{\vb{E}'}{\frac{\vb{v}}{v}}
        \frac{\vb{v}}{v}
        +
        \gamma\,\vb{v}\cross\vb{B}'
      \right)
    &=
    c^2\curl\!\left(
        \gamma\,\vb{B}'
        -
        (\gamma-1)
        \ip{\vb{B}'}{\tfrac{\vb{v}}{v}}
        \tfrac{\vb{v}}{v}
        -
        \frac{\gamma}{c^2}\vb{v}\cross\vb{E}'
      \right).
  \end{aligned}
  \label{eq:pre_final_maxwell}
\end{equation}

Hence, by the principle of relativity, the electric and magnetic fields measured in $F'$
must be

\begin{align}
  \vb{B}_{F'}
  &=
  \phi(\vb{v})\left(
      \gamma\,\vb{B}'
      -
      (\gamma-1)
      \ip{\vb{B}'}{\tfrac{\vb{v}}{v}}
      \tfrac{\vb{v}}{v}
      -
      \frac{\gamma}{c^2}\vb{v}\cross\vb{E}'
    \right),
  \\
  \vb{E}_{F'}
  &=
  \phi(\vb{v})\left(
      \gamma\,\vb{E}'
      -
      (\gamma-1)
      \ip{\vb{E}'}{\tfrac{\vb{v}}{v}}
      \tfrac{\vb{v}}{v}
      +
      \gamma\,\vb{v}\cross\vb{B}'
    \right),
\end{align}

where $\phi(\vb{v})$ is a scalar function of the boost velocity.

The possible values of $\phi(\vb{v})$ can be constrained by requiring that a boost by
$\vb{v}$ followed by the inverse boost $-\vb{v}$ reproduce the original fields:

\begin{align}
  \vb{B}
  &=
  \phi(-\vb{v})\left(
      \gamma\,\vb{B}_{F'}
      -
      (\gamma-1)
      \ip{\vb{B}_{F'}}{\tfrac{\vb{v}}{v}}
      \tfrac{\vb{v}}{v}
      +
      \frac{\gamma}{c^2}\vb{v}\cross\vb{E}_{F'}
    \right),
  \\
  \vb{E}
  &=
  \phi(-\vb{v})\left(
      \gamma\,\vb{E}_{F'}
      -
      (\gamma-1)
      \ip{\vb{E}_{F'}}{\tfrac{\vb{v}}{v}}
      \tfrac{\vb{v}}{v}
      -
      \gamma\,\vb{v}\cross\vb{B}_{F'}
    \right).
\end{align}

This gives
\[
  \phi(-\vb{v})\phi(\vb{v})=1.
\]

By spatial isotropy, the scalar $\phi$ can depend only on the magnitude of the boost
velocity, and therefore
\[
  \phi(-\vb{v})=\phi(\vb{v}).
\]
Consequently,
\[
  \phi(\vb{v})^2=1.
\]

The transformed fields must approach the original fields as $\vb{v}\to\vb{0}$:

\begin{align}
  \lim_{\vb{v}\to\vb{0}}\vb{B}_{F'} &= \vb{B},
  \\
  \lim_{\vb{v}\to\vb{0}}\phi(\vb{v}) &= 1.
\end{align}
Continuity therefore selects
\[
  \phi(\vb{v})=1.
\]

The final transformation laws for the electric and magnetic fields are
\begin{equation}
  \begin{aligned}
    \vb{B}_{F'}
    &=
    \gamma\,\vb{B}
    -
    (\gamma-1)
    \ip{\vb{B}}{\tfrac{\vb{v}}{v}}
    \tfrac{\vb{v}}{v}
    -
    \frac{\gamma}{c^2}\,
    \vb{v}\cross\vb{E},
    \\
    \vb{E}_{F'}
    &=
    \gamma\,\vb{E}
    -
    (\gamma-1)
    \ip{\vb{E}}{\tfrac{\vb{v}}{v}}
    \tfrac{\vb{v}}{v}
    -
    \gamma\,\vb{v}\cross\vb{B}.
  \end{aligned}
  \label{eq:relativistic_maxwell}
\end{equation}
 

\end{document}
